\section{Which transfers are worth pursuing next?}
\label{sec:programme}

The new distinctions change the order in which proposed extensions should
be tested. A successful analytic criterion can coexist with an impossible
Selmer instantiation. Conversely, an admissible Selmer representation need
not produce the desired prime statistic. The following elementary fact
helps distinguish the nonabelian RH and BSD directions.

\begin{proposition}[Finite-image representations pass the rational-duality test]
Let $\rho:G_K\to\GL(V)$ have finite image, and suppose the cyclotomic
character of $G_K$ has infinite image. Then
\[
 \Hom_{F[G_K]}(V,(V^c)^\vee(1))=0.
\]
For any stable lattice the full congruence term is therefore finite.
\end{proposition}
\begin{proof}
The common kernel of $\rho$ and $\rho^c$ has finite index. The cyclotomic
character still has infinite image on it. Choose an element $g$ in this
kernel with $\chi(g)\ne1$. It acts as identity on $V$ and as the scalar
$\chi(g)$ on $(V^c)^\vee(1)$. Every equivariant map is consequently killed
by the nonzero scalar $1-\chi(g)$ and must be zero. The exact finiteness
criterion applies to a stable lattice. More explicitly, at torsion level
$n$ every matrix entry is killed by $1-\chi(g)$, giving the uniform bound
$\len\Hom_{G_K}(T/\pi^n,D(T)/\pi^n)\le(\dim V)^2v_\pi(1-\chi(g))$.
For a number field the cyclotomic image is infinite, since adjoining all
$\ell$-power roots of unity has infinite degree.
\end{proof}

This removes one obstruction for finite-image Artin targets. It does not
establish residual irreducibility at every coefficient prime, the source's
scalar-inertia condition, the required local cocycle ranks, or the
arithmetic theta-divisibility identity. It also does not establish Artin
holomorphy. Its role is to show why the central polarized obstruction to
BSD should not be used to dismiss every nonabelian zero problem.

\subsection{An explicit accounting rule for any new comparison}

\begin{lemma}[Transfer budget]
Let $D\ge n\ge\theta D>0$, $m>0$, $A\ge0$, $U,L,J\ge0$, and $R,V,I\in\R$.
Suppose
\[
 D(A-L)\le V+2I,\qquad V\le n(A-mR+U),\qquad I\le DJ.
\]
Then
\[
 mR\le U+\frac{L+2J}{\theta}.
\]
\end{lemma}
\begin{proof}
Combining the assumptions and rearranging gives
\[
 nmR\le(n-D)A+nU+D(L+2J)\le nU+D(L+2J).
\]
Divide by $n>0$ and use $D/n\le1/\theta$.
\end{proof}

This makes two risks explicit. First, a very thin chosen component can
have $\theta\to0$, losing the exponent even when all other estimates are
correct. Second, a new projector or congruence loss must enter the budget.
Clearing $n$ columns by one scalar $d$ multiplies the squared Gram
determinant by $|d|^{2n}$, adding $2n\log|d|$; the equivalent local
description charges the corresponding wedge valuations. These are two
descriptions of the same clearing operation, not independent costs to
charge twice.

At the source's final normalization, an added per-dimension cost $E(x)$
would contribute on the scale
\[
 \frac{E(x)}{\theta\,m\,a\,\#\Phi_b}.
\]
Thus a sufficient scale condition is
$E(x)/(a\#\Phi_b)\ll_\varepsilon x^\varepsilon$ and
$\theta^{-1}\ll_\varepsilon x^\varepsilon$ for every $\varepsilon>0$,
with $m\asymp\sqrt{x}$. A fixed projector denominator is harmless for
this asymptotic purpose. A denominator or missing-dimension ratio growing
with the source's rapidly changing weights, levels, and ranks is not
automatically harmless. The source does not provide a fixed embedding of
an arbitrary prescribed eigenform throughout that family.

A bound for a full determinant does not give a bound for a selected
factor: $\operatorname{diag}(u,u^{-1})$ has determinant one while its first
factor tends to zero as $u\to0^+$. This elementary example is why an
individual-form lower bound needs proof. A whole-orbit approach avoids
some spectral splitting but still needs the matching arithmetic
construction.

\subsection{Three concrete research routes}

\paragraph{1. Formalize the symmetric mirror bound first.}
This is the smallest substantive extension identified here. Introduce the
symmetric conjugate-dual submodule, prove the image inclusion using
\texttt{mirrorForm\_eq\_zero}, and replace the final length comparison
by the sharper one. It can first be proved for the current
\texttt{SelmerDatum}, without altering a public RH theorem. Then split
off a raw datum that does not assume $C(T)<\infty$, and propagate a finite
$C^+(T)$ through descent and torsion specialization. The exact-sign theorem
provides a clean algebraic class that a generalized theorem could admit.

No concrete new arithmetic representation satisfying every remaining
hypothesis has been constructed here. Moreover, relevant automorphic
polarization signs are constrained by the theorem of
Bella\"iche--Chenevier~\cite{bc}; their sign convention and multiplier must
be translated before applying it to this interface. One cannot assume
that opposite-sign examples include the central geometric motives of
interest. The direct Weil-pairing calculation already proves that they
do not include the BSD Tate modules considered above.

\paragraph{2. Build one arithmetic realization of a nonabelian statistic.}
The previous reports have already reduced nonabelian zero problems to
explicit inequalities, including the cubic discriminant-$-23$ example.
The next contribution should prove a divisibility or covolume lemma that
distinguishes the genuinely nonabelian Frobenius classes, while preserving
the same period normalization, coverage, and denominator budget. A
finite-image target has no rational conjugate-Tate-duality obstruction,
which makes this direction more compatible with the present Selmer
architecture than a central polarized target. Even the first cubic case
requires new arithmetic; no such lemma is proved in this report.

\paragraph{3. Use an entire modular orbit as the target block.}
Start with the two-dimensional weight-24 space, or the one-dimensional
space generated by $\Delta$. For the first, keep both conjugate eigenforms
together; for the second there is no spectral projector cost to begin
with. Construct the local trace statistic and a compatible theta-lift
comparison. If the Selmer objects are chosen at a central self-dual
normalization, the symmetric global obstruction remains. An alternative
using the source's positive weight separation must explicitly connect
those off-central Selmer objects to the desired $L$-function; a Tate shift
cannot simply be treated as changing only notation.

\subsection{Status of candidate problems}

\begin{longtable}{@{}>{\raggedright\arraybackslash}p{0.24\textwidth}>{\raggedright\arraybackslash}p{0.31\textwidth}>{\raggedright\arraybackslash}p{0.39\textwidth}@{}}
\toprule
Target & What this analysis establishes & Remaining decisive work\\\midrule
\endhead
Nonabelian Dedekind / Artin RH & Finite-image representations pass the
rational-duality test; earlier analytic detectors remain available. & A
new Frobenius-sensitive divisibility estimate, matching periods and
denominators; holomorphy where a factor argument requires it.\\
Modular / symmetric-power RH & Exact spectral costs, whole-orbit alternative,
and inertia/self-duality screening. & Arithmetic realization of normalized
Hecke traces with compatible global lower bounds.\\
Generalized anticyclotomic Selmer selection & Proved symmetric mirror bound
and exact sign dichotomy. & Refactor and verify the full theorem; exhibit
an arithmetic instance satisfying the remaining hypotheses.\\
BSD / central polarized Bloch--Kato targets & Direct instantiation fails;
the surviving polarization component is identified. & A rank-sensitive
descent or new estimate for the mirror preimage of genuine pairings.\\
Vandiver / higher $K$-group vanishing & No new vanishing claimed. Tate
representations and untwisting still require separate admissibility and
comparison arguments. & The arithmetic comparison and exceptional-prime
conditions already isolated in the predecessors.\\
\bottomrule
\end{longtable}

The strongest deliverable here is a precise account of the arithmetic
requirements, including a proved improvement of one internal bound. It
provides a concrete next theorem to formalize and explains why that theorem
alone will not settle the central polarized cases.
